\section{Irreducible restrictions and the chain bound}
\label{sec:rigidity}

We finish the proof of Theorem~\ref{cls:chain-rigidity}.  The essential
representation-theoretic input is the corrected classification of irreducible
triples in \cite[Theorem~4.1 and Table~3]{BurnessTesterman2018}.
The composition-length bound of Section~\ref{pkt:section} is the starting
point.  After stabilizing the numbers of trivial and nontrivial factors,
a first rank bound forces the natural endpoint modules to restrict
irreducibly on earlier endpoints.  A second rank bound uses this new
irreducibility to limit the source ranks of the nonlinear edges on later
shortcuts.  The long chain supplies more such ranks than the bound permits.

We develop the two classification consequences together.  The first uses
bounded composition length of a natural restriction; the second uses an
intervening simple group of larger rank whose natural module remains
irreducible.  These hypotheses control the twisted-diagonal spin families
in different ways.  Simplicity of the smallest group alone gives no
uniform bound on the ambient ranks.

\subsection{The two rank bounds}

Throughout this section, algebraic groups are defined over an algebraically
closed field $k$ of characteristic $p>0$, and representations are rational.
We pass to simply connected covers when defining highest weights.  For a
classical simple group $J$, write $N_J$ for its natural simple module and
$n(J)=\dim N_J$.  In characteristic two, for $J$ of type $B_r$, this means the
$2r$-dimensional simple quotient of the orthogonal natural module of dimension
$2r+1$.  Thus the finite $B_r$ and $C_r$ conventions agree.  Write
$\ell_H(M)$ for the composition length of an $H$-module $M$.

We call a nontrivial simple module \emph{of natural type} if it is a Frobenius
twist of a natural simple module or its dual, allowing the central, diagram,
and characteristic-two $B/C$ isogenies in the chosen realization.  At ranks
at least nine, its highest weight on a fixed simply connected source is
$p^a\omega_1$, or $p^a\omega_r$ for the dual in type $A_r$, for some $a\geq0$.
The exclusions of natural type below therefore include the natural module,
its dual, and the simple natural quotient in the characteristic-two odd
orthogonal realization.

We record explicitly the image convention in the irreducible-triple Theorem.
If $\varphi:Y\longrightarrow\operatorname{SL}(V)$ is its representation, the
table records $(\varphi(Y),\varphi(H),V)$.  These images have the same ranks as
the original groups.  The possible $B/C$ change in characteristic two must
be made using this fixed representation; see
\cite[Lemma~2.1 and Remark~4.2(b)]{BurnessTesterman2018}.
Let $\sigma:B_r\longrightarrow C_r$ and $\tau:C_r\longrightarrow B_r$ be the
special isogenies.  With Bourbaki numbering, their pullbacks on fundamental
weights are
\begin{align}
 \sigma^*(\omega_i^C)&=\omega_i^B &&(i<r), &
 \sigma^*(\omega_r^C)&=2\omega_r^B,\notag\\
 \tau^*(\omega_i^B)&=2\omega_i^C &&(i<r), &
 \tau^*(\omega_r^B)&=\omega_r^C.
 \label{irr:bc-weights}
\end{align}
Both composites are Frobenius; these are the root-isogeny formulas of
\cite[Chapter~I, Section~3.2]{DowdSin1996}.  In particular,
\begin{equation}
 \sigma^*N_{C_r}=N_{B_r},\qquad
 \tau^*N_{B_r}=N_{C_r}^{[1]}.
 \label{irr:bc-natural}
\end{equation}
Here $[a]$ denotes the twist by $p^a$.
Consequently the length and irreducibility of a natural restriction are
preserved on passing to the image groups in the table.  Frobenius twists
preserve these properties by transporting invariant subspaces by the
entrywise field automorphism.  This argument does not require the subgroup
embedding to be defined over the prime field.

For reference, the complete part of Table~3 with a connected simple subgroup
of rank $r\geq9$ is displayed in Table~\ref{irr:high-rank-table}.
We omit restrictions on $p$ and on the weight of $V$ that only remove cases
and hence cannot enlarge the set of ranks.  The two rows with an unbounded
number of blocks retain their essential condition $p=2$.

\begin{table}[htbp]
\centering
\small
\begin{tabular}{lll}
\toprule
Rows of Table~3 & Subgroup image & Ambient image\\
\midrule
$\mathrm I_1,\mathrm I'_1$ & $C_r$ & $A_{2r-1}$\\
$\mathrm I_2,\mathrm I_3$ & $B_r$ & $A_{2r}$\\
$\mathrm I_4,\mathrm I_5$ & $D_r$ & $A_{2r-1}$\\
$\mathrm I_6$ & $A_r$ & $A_{(r^2+r-2)/2}$\\
$\mathrm I_7$ & $A_r$ & $A_{(r^2+3r)/2}$\\
$\mathrm {IV}_1,\mathrm {IV}'_1$ & $B_r$ & $D_{r+1}$\\
$\mathrm {IV}'_2$ & diagonal $B_r$ in $B_rB_r$ & $D_{2r+1}$\\
$\mathrm S_6$ & diagonal $B_r$ in $B_r^t$ & $D_{tr+1}$, $t\geq2$, $p=2$\\
$\mathrm {MR}_4$ & $D_r$ & $C_r$\\
$\mathrm {MR}_5$ & diagonal $B_r$ in $B_r^t$ & $B_{tr}$, $t\geq2$, $p=2$\\
\bottomrule
\end{tabular}
\caption{The high-rank simple-subgroup rows of the corrected irreducible-triple
table.  Diagonal projections may involve different Frobenius morphisms.}
\label{irr:high-rank-table}
\end{table}

Here and below ``diagonal'' allows the twisted projections explained in
\cite[Remark~6.1(c)]{BurnessTesterman2018}.  A simple group projecting onto
several simple blocks has an isogeny onto each block, so their ranks all equal
$r$.  The exceptional projection of an $A_1$ into a $B_2$ block cannot occur
when $r\geq9$.  Every row not displayed either has subgroup rank at most
eight or specifies an actual nonsimple product.  This proves the exhaustion
asserted in Table~\ref{irr:high-rank-table}; one cannot replace a product in
an unprimed row by an arbitrary diagonal subgroup without again testing
irreducibility.

Put
\begin{equation}
 \mathcal R(r)=
 \left\{r,r+1,2r-1,2r,2r+1,
 \frac{r^2+r-2}{2},\frac{r^2+3r}{2}\right\}.
 \label{irr:rank-set}
\end{equation}
All rows outside $\mathrm S_6$ and $\mathrm {MR}_5$ have ambient rank in this
set, which has at most seven elements.

\begin{lemma}[Rank bound with composition length]
\label{irr:rank-length}
Let $H<Y$ be connected classical simple algebraic groups with
$r=\operatorname{rk}(H)\geq9$.  Suppose $V$ is a nontrivial irreducible,
$p$-restricted, tensor-indecomposable $Y$-module which is not of natural type,
and $V|_H$ is irreducible.  If $\ell_H(N_Y)\leq c$ for an integer $c\geq1$, then
\[
 \operatorname{rk}(Y)\in
 \mathcal R(r)\cup\{tr,tr+1:2\leq t\leq c\}.
\]
In particular the number of possible ambient ranks is at most $2c+7$.
The assertion is unchanged by central covers and the isogenies in
\eqref{irr:bc-weights}.
\end{lemma}

\begin{proof}
Apply the corrected irreducible-triple Theorem to $V$ and use
Table~\ref{irr:high-rank-table}.  Only $\mathrm S_6$ and $\mathrm {MR}_5$
require an additional argument.  Each of their $t$ blocks contributes a
nontrivial constituent of the natural restriction, of dimension $2r$.
In $\mathrm {MR}_5$, passing from the odd orthogonal natural module to its
simple quotient removes one trivial constituent and leaves all $t$
nontrivial constituents.  In $\mathrm S_6$, there are also two trivial
constituents, which do not affect the lower bound $t$ on length.  Different
Frobenius twists may change the isomorphism types of these constituents, but
not their dimensions or their number counted with multiplicity.  Thus $t\leq c$.
The image conventions and \eqref{irr:bc-natural} justify the same calculation
on the original groups.  No complete-reducibility assertion is used.
\end{proof}

The length hypothesis cannot simply be dropped.  For example, in
characteristic two fix $H=B_r$ and embed it in $B_r^t<B_{tr}$ by
$g\mapsto(g,F(g),\ldots,F^{t-1}(g))$.  The spin module of $B_{tr}$ restricts
to
\[
 \bigotimes_{a=0}^{t-1}L_H(\omega_r)^{[a]},
\]
which is irreducible by Steinberg's Tensor Product Theorem.  The ambient
spin module is restricted, nonnatural, and tensor-indecomposable, whereas
$t$ is arbitrary.  This is the simple diagonal case of $\mathrm {MR}_5$.

\begin{lemma}[Rank bound with an intervening group]
\label{irr:intervening}
Let $H<X<Y$ be connected classical simple algebraic groups with
$r=\operatorname{rk}(H)\geq9$ and $s=\operatorname{rk}(X)>r$.
Suppose $N_X|_H$ is irreducible.  If $V$ is a nontrivial irreducible,
$p$-restricted, tensor-indecomposable $Y$-module, not of natural type, and
$V|_H$ is irreducible, then
\[
 \operatorname{rk}(Y)\in\mathcal R(r).
\]
The same conclusion holds when $X=Y$ and $N_Y|_H$ is irreducible.
\end{lemma}

\begin{proof}
Use the same representation
$\varphi:Y\longrightarrow\operatorname{SL}(V)$ throughout the proof.
For $(Y,H,V)$, every row of Table~\ref{irr:high-rank-table} outside
$\mathrm S_6$ and $\mathrm {MR}_5$ gives the conclusion.  Suppose one of
these two rows occurs.  Then $p=2$, the image of $H$ is $B_r$, the image of
$Y$ is $D_m$ or $B_m$, and $V$ is a half-spin or spin module, respectively.
Let $W$ be the orthogonal natural module of this fixed ambient image, on
a cover when necessary.  Every nontrivial composition factor of $W|_H$
has dimension $2r$.

Since $H<X$, irreducibility on $H$ implies irreducibility on $X$.
Apply the irreducible-triple Theorem a second time to $(Y,X,V)$, keeping
the ambient group and representation unchanged.  The theorem requires
restrictedness and tensor-indecomposability on $Y$, so it makes no such
assumption about the $X$-module $V|_X$.
If the ambient image is $B_m$, the only row with a simple subgroup of rank
at least ten is $\mathrm {MR}_5$.
If it is $D_m$ and $V$ is half-spin, the only such rows are
$\mathrm {IV}_1$, $\mathrm {IV}'_1$, $\mathrm {IV}'_2$, and $\mathrm S_6$.
In all these cases the image of $X$ is $B_s$ and every nontrivial composition
factor of $W|_X$ is a Frobenius twist of $N_X$, of dimension $2s$.

By hypothesis and \eqref{irr:bc-natural}, each such constituent remains
irreducible on $H$.  Restricting a composition series for $W|_X$ and refining
it over $H$ therefore exhibits an $H$-composition factor of dimension $2s$.
This contradicts the preceding dimension $2r$, since $s>r$.
The argument works for nonsplit natural restrictions as well as for direct
sums.

There is no extra case from $\mathrm {MR}_4$: its ambient \emph{image} is
$C_m$, whereas the image of the fixed representation here is $B_m$ or $D_m$.
One cannot change that image between the two applications.
Equations~\eqref{irr:bc-weights}--\eqref{irr:bc-natural} preserve the natural
irreducibility hypothesis on passing to these image groups, including an
inseparable isogeny.  They do not replace the fixed spin representation $V$
by a natural representation.

Finally, if $X=Y$, the irreducibility of $N_Y|_H$ directly excludes
$\mathrm S_6$ and $\mathrm {MR}_5$, since their natural restrictions have at
least two nontrivial composition factors.
\end{proof}

\subsection{From bounded length to natural irreducibility}

We now use the finite components and algebraic lifts supplied by
Proposition~\ref{lift:edge-lifts} and Lemma~\ref{lift:path-lifts}.
Fix the last endpoint of a finite retained chain and one of its socle
coordinates.  Proposition~\ref{lift:edge-lifts} realizes the earlier endpoint
components in this coordinate as nested quasisimple subgroups: the terminal
derived groups of the nested inner projections nest.  If $P_i\le P_j$ are
two such projections and $Q_j\le R_j\le P_j$ are the core and active
preimages of Lemma~\ref{cls:core-gaps}, their original pair comparison is
$(P_i\cap R_j)/Q_j$.  The two-step component argument in
Proposition~\ref{lift:edge-lifts} identifies its terminal derived subgroup
with the conjugation image of the lower component on the higher one;
the kernel is central.  Thus this common finite realization carries the
pair restrictions bounded in Theorem~\ref{pkt:bounded-length}.

Write $L_i$ for the finite component at a chain endpoint,
$\mathbf G_i$ for its simply connected algebraic source, and $N_i$ for
its natural simple module over $k$.
In the notation $N_j|_{L_i}$, take the natural representation on the
standard finite cover $U_j$ and pull it back along the finite map
$U_i\to U_j$ from \eqref{lift:finite-diagram}, or a composite of those maps.
Two such finite lifts inducing the same map to $L_j$ differ by a
homomorphism from the perfect group $U_i$ to the central kernel of
$\pi_j$, and hence are equal.  The finite pair maps therefore agree with
their composites, so these restrictions compose.
Algebraic restrictions are taken along one chosen
composed path at a time.  We make no identification of independently chosen
shortcut lifts.

\begin{lemma}[Stabilized composition factors]
\label{irr:constituent-count}
Suppose that, for every $i<j$ in a chain of quasisimple components, the
restriction $N_j|_{L_i}$ has exactly $a$ nontrivial and $b$ trivial
composition factors.  Assume $a\geq1$ and that every nontrivial simple
$L_j$-constituent restricts nontrivially to each earlier $L_i$.
For every $i<j<k$, each nontrivial composition factor of $N_k|_{L_j}$
restricts irreducibly and nontrivially to $L_i$.
\end{lemma}

\begin{proof}
Let $U_1,\ldots,U_a$ be the nontrivial $L_j$-constituents, counted with
multiplicity.  If $U_t|_{L_i}$ had only trivial composition factors, the
image of the perfect group $L_i$ would be conjugate into an upper unitriangular
group and hence would be both perfect and solvable.  It would be trivial,
contrary to the nontriviality hypothesis.  Thus, writing $a_t,b_t$ for its
nontrivial and trivial composition counts, we have $a_t\geq1$ and $b_t\geq0$.
Additivity of composition multiplicities along restriction gives
\[
 a=\sum_{t=1}^a a_t,\qquad
 b=b+\sum_{t=1}^a b_t.
\]
It follows that $a_t=1$ and $b_t=0$ for every $t$, as required.
\end{proof}

\begin{proposition}[Natural endpoint restrictions]
\label{irr:natural-endpoints}
In the remaining configurations of Proposition~\ref{grw:remaining-paths},
arbitrarily long chains have arbitrarily long subchains on which every
natural endpoint module restricts absolutely irreducibly to each earlier
endpoint.
\end{proposition}

\begin{proof}
By Theorem~\ref{pkt:bounded-length}, fix an absolute integer $c$ bounding
$\ell_{L_i}(N_j)$ for every retained comparison $i<j$.
The pair of nontrivial and trivial composition counts has at most
$(c+1)^2$ possible values.  Finite Ramsey applied to these pair colors
gives arbitrarily long subchains with constant counts $(a,b)$.
These restrictions are nontrivial, and $a\geq1$ follows because a perfect
group cannot act nontrivially with only trivial composition factors.
Moreover, a nontrivial simple representation of a quasisimple group has
central kernel.  Its restriction along the noncentral earlier component
is therefore nontrivial.  The hypotheses of
Lemma~\ref{irr:constituent-count} hold.

For each $j<k$, realize every nontrivial constituent $U$ of $N_k|_{L_j}$
as the restriction of the corresponding restricted finite-group highest
weight module of $\mathbf G_j$.  Steinberg's Tensor Product Theorem
decomposes that rational simple module into Frobenius twists of
$p$-restricted simple modules; see \cite[Theorem~41]{SteinbergNotes}.
If a restricted factor is tensor decomposable, continue factoring it
until its nontrivial factors are tensor-indecomposable.
This terminates since every factor has smaller positive dimension.
Every resulting factor is $p$-restricted: its dominant highest weight
is a summand of a $p$-restricted dominant weight.

Color a triple $i<j<k$ red if some such $U$ has a nontrivial terminal
factor $V$, not of natural type, and blue otherwise.  In a red triple,
Lemma~\ref{irr:constituent-count} says that $U|_{L_i}$ is irreducible.
Hence each tensor factor $V^{[d]}$ is irreducible on $L_i$.
Use a path lift from $\mathbf G_i$ to $\mathbf G_j$ and replace its image
by its $p^d$-Frobenius image.  The module $V$ is irreducible on this closed
connected simple subgroup: any proper rational invariant subspace would
also be invariant under the corresponding finite component.  The subgroup
has rank $r_i$ and is proper because the endpoint ranks increase.

The natural restriction on this subgroup has algebraic composition length
at most $c$.  Indeed, restriction from an algebraic group to a finite
subgroup is exact and cannot reduce composition length, and Frobenius
transport preserves length.  The finite restriction here is a twist of
$N_j|_{L_i}$.
Lemma~\ref{irr:rank-length} consequently bounds the number of possible
middle ranks $r_j$ by $2c+7$, for fixed $i$.

A red homogeneous chain with more than $2c+9$ vertices is impossible:
fix its first and last vertices $i,k$ and vary its middle vertex $j$.
The distinct ranks $r_j$ would exceed that bound.
A second application of finite Ramsey therefore gives arbitrarily long
blue homogeneous subchains.

Fix $i<j<k$ in a blue subchain and choose a nontrivial constituent $U$
of $N_k|_{L_j}$.  There is at least one nontrivial terminal tensor factor,
and every such factor is now of natural type.
Since $U|_{L_i}$ is irreducible, a natural or dual natural Frobenius factor
restricts irreducibly to $L_i$.
Undoing its twist, dual, and the isogenies
\eqref{irr:bc-weights} shows that $N_j|_{L_i}$ is irreducible.
Delete the last vertex of the blue subchain, so that a suitable $k$ exists
for every remaining pair $i<j$.
All modules and restrictions were taken over $k$, so the resulting finite
restrictions are absolutely irreducible.
\end{proof}

\subsection{Restricted natural factors and shortcut tails}

The word ``restricted'' in the next lemma concerns one fixed algebraic
source.  A representation with many different Steinberg digits need not
have bounded tensor length.  The retained nonlinear edges of
Proposition~\ref{grw:remaining-paths}, however, have one digit after a
common Frobenius twist is removed.

\begin{lemma}[A restricted product of natural factors]
\label{irr:restricted-natural}
Let $J$ be a classical simple algebraic group of rank at least nine.
Let $M$ be a nontrivial $p$-restricted simple $J$-module.
If every nontrivial tensor-indecomposable factor of $M$ is of natural type,
then $M\cong N_J$ or $N_J^*$.
In particular there is exactly one nontrivial terminal factor and
$\dim M=n(J)$.
\end{lemma}

\begin{proof}
A tensor factor of a simple module is simple, since any proper nonzero
submodule of the factor would tensor to one of the whole module.
Highest weights add in an irreducible tensor product.
The highest weights of the natural-type factors have the form
$p^a\omega_1$, or $p^a\omega_r$ in type $A_r$ for the dual.
Equation~\eqref{irr:bc-weights} shows explicitly that the special isogenies
add no further natural weights.
Since their sum is $p$-restricted and all fundamental-weight coefficients
are nonnegative, every exponent $a$ is zero.
The terminal factors are therefore $N_J$ or $N_J^*$.

Any two of these factors have reducible tensor product.
For two identical factors, the kernel of the flip minus the identity is
a nonzero proper invariant subspace: it contains $v\otimes v$, while
$v\otimes w-w\otimes v\ne0$ for independent $v,w$.
This also proves properness in characteristic two.
For opposite factors, $N_J\otimes N_J^*\cong\operatorname{End}(N_J)$
contains the proper invariant line of scalar endomorphisms.
The same arguments apply to dual identical factors.
Tensoring such a proper submodule with all remaining factors contradicts
simplicity of $M$.  There is thus exactly one terminal factor.
\end{proof}

\begin{lemma}[Pulling irreducibility through a natural tail]
\label{irr:natural-tail}
Consider an edge between the finite components associated with algebraic
sources $J$ and $K$.  Suppose every nontrivial absolute composition factor
of the target natural module, restricted to the source, is a tensor product
of natural-type $J$-modules.  If that target module is absolutely irreducible
on an earlier noncentral quasisimple component $L$, then $N_J|_L$ is
absolutely irreducible.
This applies to each retained geometric edge and to every one-digit edge
whose restricted terminal factors are all of natural type.
\end{lemma}

\begin{proof}
Irreducibility on $L$ implies irreducibility on the source component, so
the target restriction has a single absolute composition factor.
It is nontrivial because the target natural dimension exceeds one.
By hypothesis this factor is a tensor product of natural-type modules.
Irreducibility of its further restriction forces each factor to restrict
irreducibly.  Undoing a twist, a dual, and the natural isogenies yields the
assertion for $N_J$.
For retained geometric edges this also follows from
Corollary~\ref{lift:geometric-irreducibility}.  For the stated one-digit edges it follows
from Lemma~\ref{irr:restricted-natural}, after restoring the common twist.
\end{proof}

\subsection{The final contradiction}

\begin{proof}[Completion of the proof of Theorem~\ref{cls:chain-rigidity}]
Suppose tested chains with different endpoint labels have unbounded length.
Apply Proposition~\ref{cls:reduction} and
Proposition~\ref{grw:remaining-paths}.
We obtain arbitrarily long configurations with strictly increasing endpoint
ranks, with their minimum rank tending to infinity, and with two-edge
shortcuts on all pairs.  There are absolute constants $c_0>0$, $C\geq1$,
and $B\geq1$ such that the natural endpoint dimensions satisfy
\begin{equation}
 n_j\geq c_0 n_i^2\qquad(i<j),
 \label{irr:quadratic-growth}
\end{equation}
and every retained geometric edge satisfies
\begin{equation}
 n_{\mathrm{target}}\leq C n_{\mathrm{source}}^B.
 \label{irr:geometric-degree}
\end{equation}
Every other edge is a defining-characteristic irreducible representation
with one restricted digit after removing a common twist.  At least one
such nonlinear edge occurs in each shortcut.
Proposition~\ref{lift:edge-lifts} and Lemma~\ref{lift:path-lifts} supply the actual
component homomorphisms and their algebraic path lifts.
By Proposition~\ref{irr:natural-endpoints}, pass to arbitrarily long
subchains with all natural endpoint restrictions absolutely irreducible.

Color a pair of endpoints red if its chosen shortcut has a one-digit edge
with a nontrivial restricted tensor-indecomposable factor not of natural
type.  Color the pair blue otherwise.
Finite Ramsey leaves arbitrarily long homogeneous configurations of one
of these two kinds.  We exclude both.

First consider a blue configuration.  On every one-digit edge,
Lemma~\ref{irr:restricted-natural} bounds the target natural dimension by
the source natural dimension; restoring a Frobenius twist does not change
dimension.  Thus every edge of every shortcut satisfies
\eqref{irr:geometric-degree}, enlarging $C,B$ if necessary.
Composing at most two edges gives the uniform bound
\begin{equation}
 n_j\leq A n_i^E\qquad(i<j),\qquad
 A=C^{B+1},\quad E=B^2.
 \label{irr:pair-degree}
\end{equation}
Choose a fixed integer $d$ with $2^d>E$.
Along $d$ consecutive comparisons of a retained chain,
\eqref{irr:quadratic-growth} gives
\[
 n_d\geq c_0^{\,2^d-1}n_0^{2^d},
\]
whereas the shortcut from its first to its last endpoint gives
$n_d\leq A n_0^E$.
These inequalities contradict each other as the minimum rank, and hence
$n_0$, tends to infinity.  This proves the bounded-natural-factor case
without imposing a bound on the number of Steinberg digits of an arbitrary
composite endpoint representation.

Now consider a red homogeneous configuration.  Fix an early endpoint $H$
of rank $r\geq9$.  Choose eight later pair comparisons with disjoint ordered
index intervals, leaving a gap between successive intervals.
For each such comparison, let $X$ be its lower endpoint, and choose the
last edge in its shortcut that has a restricted factor not of natural
type.  Let $Y$ be the algebraic source of that edge and $Z$ its target.
Every edge after it has a natural tail in the sense of
Lemma~\ref{irr:natural-tail}: it is either geometric or a one-digit edge
with only natural restricted terminal factors.
The upper endpoint natural module is irreducible on the finite component
of $H$.  Applying that lemma backwards along the tail shows that $N_Z$
is irreducible on the same component.

Remove the common Frobenius twist of the $Y$-representation on $N_Z$,
and choose a nontrivial restricted tensor-indecomposable factor $V$ not
of natural type.  Irreducibility of the whole restriction implies
irreducibility of the corresponding twisted factor on the finite $H$.
Compose the actual algebraic maps along the path from $H$ through $X$
to $Y$, and apply the removed Frobenius morphism to both lower images.
Then $V$ is irreducible on the resulting connected simple image of $H$.
Its rank is $r$, and the resulting image of $X$ has its original rank
$s>r$.  Its natural simple module remains irreducible on the image of
$H$, by the endpoint conclusion and \eqref{irr:bc-natural}.
Finite irreducibility implies each of these rational irreducibility
statements, since a rational invariant subspace would also be invariant
under the actual finite image.

If the algebraic image of $X$ is proper in $Y$, apply
Lemma~\ref{irr:intervening}.  If that image is all of $Y$, the map from
$X$ to $Y$ is an isogeny.  This can occur even when the finite groups are
strictly nested, for example on a subfield edge.  The pullback of $N_Y$
is a natural simple module or its dual, up to Frobenius twist and the
isogenies in \eqref{irr:bc-natural}, by the root-datum classification
of isogenies
\cite[Definition~6.1.8, Example~6.1.9 and Theorem~6.1.16(1)]{Conrad2014}.
Indeed, for compatible simple roots an isogeny has a root bijection
$\delta$ and powers $q_i$ of the characteristic.  Write the Cartan
matrices in the row-coroot convention
$A^X_{ij}=\langle\alpha_j,\alpha_i^\vee\rangle$ and
$A^Y_{ab}=\langle\beta_b,\beta_a^\vee\rangle$.
The compatibility equation is
$q_j A^X_{ij}=q_i A^Y_{\delta(i)\delta(j)}$.
Pairing the pullback of the natural fundamental weight with the source
simple coroots gives $q_i$ at the node mapped to the natural node and
zero elsewhere.  In the retained ranks this node is the natural node,
or its type-$A$ dual; the only change of root lengths is the
characteristic-two $B/C$ case in \eqref{irr:bc-natural}.
Since $N_X$ is irreducible on the
lower $H$, it follows that $N_Y$ is irreducible on the corresponding
image of $H$.  The $X=Y$ case of Lemma~\ref{irr:intervening} therefore
applies.  Thus every one of the eight chosen source ranks belongs to
$\mathcal R(r)$.
They are nevertheless distinct.  Along an algebraic edge the rank cannot
decrease, because a maximal torus maps with finite kernel into a torus
of the image.  Hence each chosen source rank lies between the ranks of
the two endpoints of its comparison.  The gaps between the ordered
comparison intervals, and strict increase of endpoint ranks, make these
eight intervals disjoint in rank.
This contradicts $|\mathcal R(r)|\leq7$.

Both homogeneous alternatives are impossible.  All Ramsey bounds and
all path lengths used here depend only on the absolute constants in the
preceding lemmas.  In particular only fixed finite paths are needed
before taking the large-rank limit.  The assumed unbounded family of
tested chains cannot exist, proving the asserted absolute chain bound
and completing Theorem~\ref{cls:chain-rigidity}.
\end{proof}
